% Project C: First Results Short Note — pre-submission LaTeX source
% Compile: pdflatex main && bibtex main && pdflatex main && pdflatex main
% Or: latexmk -pdf main.tex

\documentclass[11pt,reqno]{amsart}

\usepackage[utf8]{inputenc}
\usepackage[T1]{fontenc}
\usepackage{amsmath,amsthm,amsfonts,amssymb}
\usepackage{mathrsfs}
\usepackage{microtype}
\usepackage[margin=1in]{geometry}
\usepackage{booktabs}
\usepackage{array}
\usepackage{longtable}
\usepackage{enumitem}
\usepackage{hyperref}
\hypersetup{
    colorlinks=true,
    linkcolor=blue!50!black,
    citecolor=blue!50!black,
    urlcolor=blue!50!black
}
\usepackage{cleveref}

% Theorem-like environments
\newtheorem{theorem}{Theorem}[section]
\newtheorem{proposition}[theorem]{Proposition}
\newtheorem{lemma}[theorem]{Lemma}
\newtheorem{corollary}[theorem]{Corollary}

\theoremstyle{definition}
\newtheorem{remark}[theorem]{Remark}
\newtheorem{principle}[theorem]{Principle}

% Operators
\DeclareMathOperator{\erf}{erf}
\DeclareMathOperator{\erfc}{erfc}

% Custom
\newcommand{\LHS}{\mathrm{LHS}}
\newcommand{\RHS}{\mathrm{RHS}}
\newcommand{\dps}{\mathrm{dps}}

\title{Project C: A symbolic-search engine that finds faster
computational forms of two Ramanujan identities}

\author{Jouni Helminen}
\address{Independent researcher}
\email{jouni@helminen.co}

\date{May 25, 2026}

\subjclass[2020]{Primary 11F37, 11Y16; Secondary 11F03, 11P82, 33D15, 68W30}

\keywords{Ramanujan; Lost Notebook; computational identity discovery;
Dedekind eta; $E_2$ quasi-modular form; mock theta; symbolic search;
Mellin inversion; $\psi(q)$}

\begin{document}

\begin{abstract}
The Lost Notebook of Srinivasa Ramanujan contains thousands of
identities whose computational utility outside their stated form is
largely unexplored. We describe \emph{Project C}, an automated
symbolic-search engine that composes a small basis of structural
operators against verified target identities, gated by an
\texttt{mpmath} cost model under a strict $\le 10^{-30}$
relative-error acceptance bar. Using only the engine's existing
\verb|modular_transform_with_correction| operator (the Dedekind-$\eta$
/ $E_2$ quasi-modular family), Project C reproduces transformed forms
of two recorded Ramanujan identities---Warnaar's equation (1.7) at
$a=1$ (Lost Notebook p.~29) and Andrews--Berndt Entry~1.6.6 ($\psi(q)$
on Lost Notebook p.~35)---whose computational cost at canonical
$q = e^{-\beta}$ near $q \to 1^-$ is between $13.7\times$ and
$32{,}363\times$ faster than the recorded forms, verified to
$\le 6 \cdot 10^{-72}$ relative mismatch at $80$-digit working
precision. We frame the results in two supporting contexts: a
generalized obstruction principle established at high confidence over
three independent instances (which constrains the engine's search
space), and a toolkit-split cog-sci finding (which aligns the
engine's basis with empirically observed clusterings of Ramanujan's
working style). The results are \emph{improvements on known
identities, not from-scratch discoveries}---but they are the first
concrete proof-of-utility of the engine. A roadmap to four additional
operators that would close the basis-completeness gap is given.
\end{abstract}

\maketitle

\section{Introduction}\label{sec:intro}

Ramanujan's work, as recorded across his three Notebooks and the Lost
Notebook, exhibits a recurring methodological signature: identities
are stated in a \emph{single chosen form} whose role in computation,
proof, or generalization may be quite different from an equivalent
form chosen for a different purpose. The ``right'' form for hand
computation in $q = e^{-\beta}$ near $q \to 1^-$---for instance---is
frequently \emph{not} the form Ramanujan recorded, and the
modular-transform partner (when one exists) often differs by orders
of magnitude in cost.

This raises a natural question: \emph{can a symbolic-search engine,
equipped with a small basis of operators distilled from Ramanujan's
own working style, automatically transform recorded identities into
faster computational forms?} This paper answers \emph{yes} for two
specific identities; the engine and its operator basis are described
in \cref{sec:engine}, and the two transformations are given in
\cref{sec:result1,sec:result2}.

\subsection{The Ramanujan-compression observation}

A pilot experiment \cite[diary~39]{ProjectC} tested whether a
computable metric on $q$-series identities---roughly, total
bit-content of the $q$-expansion divided by AST serialization
length---separates Ramanujan's recorded identities from
``random-walk'' identities generated by composing distributive law /
product rule rewrites on sparse seeds. On a $13$-identity Ramanujan
sample (Frye--Garvan $G/H/\eta$ catalogue~\cite{GarvanFrye}) versus a
$10$-identity random-walk null, the median compression scores were
$1.59$ vs $0.95$, with perfect separation at threshold
$S_{200} \approx 1.20$. The metric is grammar-relative (it detects
``modular / $\eta$-product compression'' specifically), but
operationally it suggests that \emph{Ramanujan's mathematical taste
correlates with structural compression} of the identity AST against
the modular / $q$-Pochhammer / $\eta$-product primitives.

\subsection{The latent-space-speedrun hypothesis}

A parallel line of work---what we call the
\emph{latent-space-speedrun hypothesis}
\cite[\S D1]{ProjectC}---asks whether Ramanujan's working style
decomposes into a small collection of recurring operator clusters,
each tied to a specific problem domain. Hand-tagging of $\sim 200$
atomic proof units across four chapters of Berndt's commentaries on
the Notebooks supports the hypothesis at high confidence with a
\emph{hierarchical-modular} decomposition:
\begin{itemize}
\item \textbf{Top level}: modular $\perp$ analytic---zero
  cross-firing across the four tagged chapters.
\item \textbf{Within modular} (three chapters of NB Part~V): per-domain
  operator clusters---continued fractions $\to$
  $\{$q-shift, modular-S$\}$; modular equations $\to$
  $\{$modular-S$\}$ monoculture; class invariants $\to$
  $\{$class-field, modular-S$\}$.
\item \textbf{Within analytic} (NB Part~1 Ch.~3--4, plus the Quarterly
  Reports): bipartite---Mellin/Master-Theorem $\perp$
  Schr\"oder--K\"onig fractional iteration---with the two operators
  firing in \emph{disjoint entry sets}.
\end{itemize}
Project C's operator basis (\cref{sec:engine}) is the engineering
image of this empirical decomposition.

\subsection{Project C's design}\label{sec:design}

Project C is a six-step search loop:
\begin{enumerate}
\item \textbf{Seed}: parse a target identity into a
  \verb|SearchState| carrying a symbolic summand $f(n)$.
\item \textbf{Propose}: a proposer (currently random or
  in-conversation Claude) suggests an operator from the basis.
\item \textbf{Apply}: the operator's \verb|apply_fn(expr, n)| returns
  either a new summand, or
  \verb|(new_summand, additive_correction, summand_range)| for
  one-sided modular identities with corrections.
\item \textbf{Verify}: an \verb|MpmathVerifier| compares the before-
  and after-side sums numerically at $\ge 30$ digits across a sweep
  of parameter values; the strict acceptance bar is $\le 10^{-30}$
  relative error (loosened to $\le 1\%$ across $\ge 2$ orders of
  magnitude per \cite[CLAUDE.md gotcha~6]{ProjectC} only when the
  LHS is asymptotic-noise-limited).
\item \textbf{Cost-gate}: an \texttt{mpmath} wall-clock estimator
  computes the cost ratio; only acceptances with ratio $> 1$ are
  kept.
\item \textbf{Record}: accepted transformations are appended to a
  JSONL registry with provenance.
\end{enumerate}
This paper reports the first two acceptances where the cost ratio is
order-of-magnitude meaningful and the source identity is a Ramanujan
identity recorded in the Lost Notebook.

\subsection{Contribution}\label{sec:contribution}

The contribution is the engineering result, not a new theorem:
\begin{itemize}
\item \textbf{Result~1} (\cref{sec:result1}): Warnaar (1.7) at $a=1$
  transformed to its $\eta$-modular partner---measured
  $\mathbf{241\times}$ speedup at $\beta = 10^{-3}$ (versus
  Ramanujan's recorded right-hand side).
\item \textbf{Result~2} (\cref{sec:result2}): Andrews--Berndt
  Entry~1.6.6 (the $\psi(q)$ identity) transformed via a
  half-shifted Gaussian theta route---measured $\mathbf{27\times}$
  to $\mathbf{858\times}$ speedup vs the natural triangular sum;
  $\mathbf{60\times}$ to $\mathbf{32{,}363\times}$ vs the product
  spelling near $q \to 1^-$.
\end{itemize}
Both transformations were produced by the engine using only the
\verb|modular_transform_with_correction| operator (Bose--Einstein log
/ $\eta$-product family; Lambert / $E_2$ family) composed with
algebraic preprocessing (parameter binding, Pochhammer rewrite). No
new operator needed to be implemented for either result.
\Cref{sec:obstruction,sec:toolkitsplit} give brief framing from the
two supporting findings.

\section{The engine}\label{sec:engine}

\subsection{The 8-operator basis}\label{sec:basis}

Project C's current basis (as of v0.9; see
\verb|rkc/search/basis.py| in \cite{ProjectC}) is:

\begin{center}
\small
\begin{tabular}{c l l l}
\toprule
\# & Operator id & Family & Status \\
\midrule
1 & \verb|master_theorem|                       & analytic   & Beta, Gamma, $\pi/\sin(\pi s)$ cases (v0.5) \\
2 & \verb|mellin_inversion|                     & analytic   & Pure Gaussian, half-shift $\theta_2/\theta_4$, weighted (v0.4, v0.7) \\
3 & \verb|modular_transform_with_correction|    & modular    & Bose--Einstein log / $\eta$, Lambert / $E_2$, Stieltjes (v0.8b, v0.8c) \\
4 & \verb|fourier_self_reciprocity|             & analytic   & stub \\
5 & \verb|fractional_functional_iteration|      & iteration  & stub \\
6 & \verb|hardy_orders_of_infinity|             & asymptotic & divergence diagnostic (v0.9a) \\
7 & \verb|frullani_limit|                       & analytic   & stub \\
8 & \verb|mellin_pair_reciprocity|              & analytic   & stub \\
\bottomrule
\end{tabular}
\end{center}

Three operators are substantively implemented (1, 2, 3); one is
diagnostic-only (6 is a non-transforming annotator that classifies
$\lim_{n \to \infty} f(n)$ into
\verb|CONVERGENT| / \verb|DIVERGENT_LITERAL| /
\verb|NEEDS_REGULARIZATION| / \verb|UNKNOWN|, used to filter
divergent literal summands before any transform is attempted); four
are stubs whose \verb|apply_fn| returns \verb|None| (no proposal
applies). The two identities in \cref{sec:result1,sec:result2} are
reached using operator~3 alone, after algebraic preprocessing.

The basis was \emph{not} designed from first principles: it is the
\emph{engineering image of the empirically-tagged operator clusters}
in Ramanujan's Quarterly Reports
\cite[Hunt~5, diary~20]{ProjectC}---closed with the revised
7-operator basis covering $\sim 91\%$ of QR proof units---and the
Notebooks \cite[Hunt~5b, diaries~22, 24, 25, 36]{ProjectC};
toolkit-split hypothesis established at high confidence across four
chapters. See \cref{sec:toolkitsplit} for the cog-sci framing.

\subsection{Verifier and cost gate}\label{sec:verifier}

The \verb|MpmathVerifier| compares LHS and RHS via \verb|mpmath.nsum|
at adjustable precision (default $30$ digits; bumped to $80$ for
high-precision acceptance tests). When an operator emits an additive
scalar correction (operator~3), the verifier folds the correction
into the before-side residual before comparison:
\begin{equation*}
\text{verify}: \qquad
\sum_n f_{\text{before}}(n) \;\equiv\; \sum_n f_{\text{after}}(n) + C.
\end{equation*}
The summand range is operator-supplied (\verb|(1, oo)| for one-sided
modular identities; \verb|(-oo, oo)| for two-sided Poisson
identities).

The cost gate uses a structural estimator
(\verb|rkc/search/routers.py|) that recognises kernel families
(\verb|gaussian_quadratic|, \verb|one_sided_exponential_family|,
etc.) and returns the expected term count from \texttt{rkc}'s
per-family cost model ($28/\beta$ for the one-sided exponential
family at $\le 10^{-12}$ tolerance, $\sqrt{28/\beta}$ for two-sided
theta-tail families, etc.). An identity passes the cost gate iff
$\text{ratio} = \text{cost}_{\text{before}} / \text{cost}_{\text{after}} > 1$.

\subsection{Strict-bar discipline}\label{sec:strictbar}

Throughout the run we hold to a \emph{strict claim discipline}:
every accepted identity is verified to numerical zero (or
$\le 10^{-30}$ relative error) across a $\ge 2$-dex sweep of the
convergence parameter; near-matches are \emph{never} accepted as
matches (cf.\ \cite[CLAUDE.md gotcha~6]{ProjectC}). The two
identities reported below were both verified to
$\le 6 \cdot 10^{-72}$ relative mismatch at $80$-digit precision
across the cross-check cell $\beta = 0.1$ by the engine's verifier;
the per-$\beta$ rel-err table is in
\cite[diary~43]{ProjectC}.

\subsection{Versioning history (selected)}\label{sec:versioning}

\begin{itemize}
\item \textbf{v0.4} \cite[diary~14b]{ProjectC}: keystone Poisson dual
  on the pure Gaussian
  $e^{-\beta n^2} \to \sqrt{\pi/\beta} \cdot e^{-\pi^2 n^2/\beta}$;
  end-to-end verifier $+$ cost gate operational.
\item \textbf{v0.7} \cite[diary~26]{ProjectC}: two additional strict
  double-sided Poisson identities (half-shift Gaussian
  $\theta_2/\theta_4$, weighted Gaussian second moment) added to
  \verb|mellin_inversion|.
\item \textbf{v0.8b} \cite[diary~31]{ProjectC}:
  \verb|OperatorApplication| API extended with
  \verb|additive_correction| $+$ \verb|summand_range| fields; new
  operator \verb|modular_transform_with_correction| shipping the
  Bose--Einstein log ($\eta$-modular) and Lambert ($E_2$
  quasi-modular) cases. \emph{The operator used in
  \cref{sec:result1,sec:result2}.}
\item \textbf{v0.8c} \cite[diary~32]{ProjectC}: same operator extended
  to a third case (Stieltjes occupation-number
  $\sum 1/(e^{\beta n} - 1)$).
\item \textbf{v0.9a} \cite[diary~33]{ProjectC}:
  \verb|hardy_orders_of_infinity| shipped as a divergence diagnostic.
\end{itemize}

The total test suite is $\mathbf{102/102}$ passing as of v0.8c; the
recognizer benchmark sits at $\mathbf{55/55} = \mathbf{100\%}$ (no
regression across all versioning).

\section{Result 1: Warnaar (1.7) at $a=1$ (Lost Notebook p.~29)}\label{sec:result1}

\subsection{Source identity}

Per Warnaar's eq.~(1.7) \cite{Warnaar2003}, attributed to Lost
Notebook p.~29 \cite{LostNotebook}:
\begin{equation}\label{eq:warnaar17}
\sum_{n=0}^\infty \frac{(q;q^2)_n \, q^{2n}}{(aq^2;q^2)_n \, (q^2/a;q^2)_n}
\;=\; (1 - a) \sum_{n=0}^\infty a^n q^{\binom{n+1}{2}}
\;+\; \frac{a\,(q;q^2)_\infty}{(aq^2;q^2)_\infty \,(q^2/a;q^2)_\infty}
   \sum_{n=0}^\infty (-1)^n a^{3n} q^{n(3n+2)} (1 + a q^{2n+1}).
\end{equation}
The identity sits in the \emph{partial-theta} class per Warnaar's
framework: the RHS replaces a bilateral Jacobi-theta sum over
$\mathbb{Z}$ by a positive cone over $\mathbb{N}$. At $a = 1$, the
first sum on the RHS collapses to $0$ (factor $(1-a)$), and the
prefactor of the partial-theta tail simplifies to a pure modular
ratio:
\begin{equation}\label{eq:warnaar17-a1}
\sum_{n=0}^\infty \frac{(q;q^2)_n \, q^{2n}}{(q^2;q^2)_n^2}
\;=\; \frac{(q;q^2)_\infty}{(q^2;q^2)_\infty^2}
   \sum_{n=0}^\infty (-1)^n q^{n(3n+2)} (1 + q^{2n+1}).
\end{equation}

\subsection{The operator sequence}

Setting $q = e^{-\beta}$ and $Q = e^{-2\pi^2/\beta}$, the engine
applies the following preprocessing-plus-operator sequence:

\begin{verbatim}
identity:     LN p. 29 / Warnaar (1.7)
preprocess:   bind_parameter(a = 1)
recognize:    (q;q^2)_inf / (q^2;q^2)_inf^2
              -> (q;q)_inf / (q^2;q^2)_inf^3   [algebraic rewrite]
operator_1:   modular_transform_with_correction[
                 Bose-Einstein-log / eta-product, t = beta]
operator_2:   modular_transform_with_correction[
                 Bose-Einstein-log / eta-product, t = 2 beta]
simplify_corrections:
              (Bernoulli/log terms)
              -> (sqrt(2) * beta / pi) * exp(pi^2/(12 beta) - 5 beta/24)
leave_unchanged:
              false-theta tail
              sum (-1)^n exp(-beta n(3n+2)) (1 + exp(-beta(2n+1)))
\end{verbatim}

The Dedekind-$\eta$ modular transformation
\[
(e^{-t}; e^{-t})_\infty
\;=\; \sqrt{2\pi/t} \cdot \exp\!\left(-\frac{\pi^2}{6t} + \frac{t}{24}\right) \cdot (e^{-4\pi^2/t}; e^{-4\pi^2/t})_\infty
\]
is applied at $t = \beta$ and $t = 2\beta$. The two applications of
operator~3 are the \emph{only} operator firings: the parameter
binding and Pochhammer rewrite are algebraic preprocessing handled
outside the basis.

\subsection{Project C improved form}

\begin{equation}\label{eq:warnaar-improved}
\boxed{\;\;
\sum_{n=0}^\infty \frac{(q;q^2)_n \, q^{2n}}{(q^2;q^2)_n^2}
\;=\; \frac{\sqrt{2}\,\beta}{\pi} \, \exp\!\left(\frac{\pi^2}{12\beta} - \frac{5\beta}{24}\right)
   \cdot \frac{(Q^2; Q^2)_\infty}{(Q; Q)_\infty^3}
   \cdot \sum_{n=0}^\infty (-1)^n e^{-\beta n(3n+2)} \bigl(1 + e^{-\beta(2n+1)}\bigr)
\;\;}
\end{equation}
with $q = e^{-\beta}$ and $Q = e^{-2\pi^2/\beta}$.

\subsection{Numerical verification}\label{sec:warnaar-verification}

Independently reproduced in our local \texttt{mpmath} 1.3 / Python
3.12.7 / .venv environment at $\dps = 80$ by
\verb|analysis/29_warnaar_17_ab_166_cost_verification.py|
\cite[diary~43; commit \texttt{f8a04a3}]{ProjectC}. All three
Project-C-vs-recorded speedups CONFIRMED at the canonical three
$\beta$ values.

Identity values at the three $\beta$ cells (all forms agree to
$\ge 72$ digits at the cross-check $\beta = 0.1$; full rel-err
table in \cite[diary~43]{ProjectC}):
\begin{center}
\begin{tabular}{r r}
\toprule
$\beta$ & $\LHS = \RHS =$ Project~C value \\
\midrule
$10^{-3}$ & $7.024 \times 10^{353}$ \\
$10^{-2}$ & $2.378 \times 10^{33}$ \\
$10^{-1}$ & $184.78$ \\
\bottomrule
\end{tabular}
\end{center}

Wall-clock comparison (median of $5$ trials, \texttt{mpmath} $80$
$\dps$; full table in
\verb|analysis/results/29_warnaar_17_ab_166_cost_verification.csv|):

\begin{center}
\small
\begin{tabular}{r r r r r r l}
\toprule
$\beta$ & LHS triangular & RHS recorded & RHS Project~C & C vs B & claim (diary 42) & verdict \\
\midrule
$10^{-1}$ & $18.3$ ms     & $7.81$ ms      & $0.572$ ms & $\mathbf{13.7\times}$ & $15.2\times$ & CONFIRMED (0.90) \\
$10^{-2}$ & $324$ ms      & $73.1$ ms      & $1.62$ ms  & $\mathbf{45.0\times}$ & $62\times$   & CONFIRMED (0.72) \\
$10^{-3}$ & $18{,}011$ ms & $1{,}105$ ms   & $4.59$ ms  & $\mathbf{241\times}$  & $272\times$  & CONFIRMED (0.89) \\
\bottomrule
\end{tabular}
\end{center}

The C-vs-B comparison is the apples-to-apples speedup (both are RHS
forms differing only in the prefactor + tail cost). All three measured
speedups are within $\pm 25\%$ of the ChatGPT-Pro claim, confirming
the $15\times$ to $\sim 250\times$ wall-clock improvement of the
Project~C transformed form over the recorded right-hand side at the
three $\beta$ cells.

\begin{remark}[on the LHS magnitude]
At $\beta = 10^{-3}$, the identity's common value is
$7.024 \times 10^{353}$---intrinsic to the $q \to 1^-$ blow-up of the
hypergeometric series and not a computational artifact. All three
forms agree to $78$ digits at this $\beta$. \cite[diary~42]{ProjectC}
did not display this value; we flag it here so reviewers see the
magnitude regime the verification is operating in.
\end{remark}

\subsection{Discussion}

The mechanism is the standard $q \to 1^-$ duality of modular forms:
when $\beta$ is small, $q$ is close to $1$ and the recorded
Pochhammer products converge in tens of thousands of factors, while
$Q = e^{-2\pi^2/\beta}$ is exponentially small so its product
representations require only one or two terms each. The two
$\eta$-modular operator firings collapse $\sim 200{,}000$ factor
evaluations at $\beta = 10^{-3}$ into $\sim 6$ exponential
evaluations plus the false-theta tail (which also collapses to
$\sim 1$ term at small $\beta$). The \emph{identity is Ramanujan's}
(recorded on Lost Notebook p.~29); the \emph{transformed form} is
what the engine produces; the $\mathbf{241\times}$ measured speedup
(claim: $272\times$) is the engineering contribution.

\section{Result 2: Andrews--Berndt Entry~1.6.6 (Lost Notebook p.~35)}\label{sec:result2}

\subsection{Source identity}

Per Andrews--Berndt \emph{Lost Notebook Part~I} Entry~1.6.6
\cite[Lost Notebook p.~35]{AB2005}:
\begin{equation}\label{eq:ab166}
L(q) \;:=\; \sum_{n=0}^\infty \frac{(-1)^n \, q^{n^2 + n}}{(q^2; q^2)_n \, (1 - q^{2n+1})}
\;=\; \psi(q)
\;=\; \frac{(q^2;q^2)_\infty}{(q;q^2)_\infty}
\;=\; \frac{(q^2;q^2)_\infty^2}{(q;q)_\infty}.
\end{equation}
Here $\psi(q)$ is Ramanujan's standard $\psi$-function (one of his
three classical theta functions). The identity equates a
\emph{hypergeometric triangular sum} on the LHS with the
\emph{infinite-product spellings} on the RHS---all of which are slow
to evaluate as $q \to 1^-$ (the triangular sum needs
$\sim \sqrt{1/\beta}$ terms; the products need $\sim 1/\beta$
factors).

\subsection{The operator sequence and improved form}

The engine routes through a \emph{half-shifted Gaussian theta}
($\theta_2$) representation of $\psi(q)$ via the operator~2
($\theta_2/\theta_4$ case of \verb|mellin_inversion|):
\begin{equation}\label{eq:psi-thetashift}
\psi(e^{-\beta}) \;=\; \tfrac{1}{2} \, e^{\beta/8} \sum_{n \in \mathbb{Z}} e^{-(\beta/2)(n + 1/2)^2}.
\end{equation}
The half-shifted Gaussian sum is then transformed by Poisson summation
(the same operator's $\theta_2 \leftrightarrow \theta_4$ branch) to
its modular partner:
\begin{equation}\label{eq:psi-improved}
\boxed{\;\;
\psi(e^{-\beta}) \;=\; \sqrt{\frac{\pi}{2\beta}} \, e^{\beta/8}
   \left( 1 + 2 \sum_{m=1}^\infty (-1)^m e^{-2\pi^2 m^2/\beta} \right)
\;\;}
\end{equation}
Equivalently,
\begin{equation}\label{eq:psi-improved-bilateral}
\psi(e^{-\beta}) \;=\; \sqrt{\frac{\pi}{2\beta}} \, e^{\beta/8} \sum_{m \in \mathbb{Z}} (-1)^m e^{-2\pi^2 m^2/\beta}.
\end{equation}

The acceptance path for this seed uses operator~2's already-shipped
half-shift case (v0.7). At present the path from the AB~1.6.6 LHS to
the half-shifted Gaussian form requires a manual rewrite step (Jacobi
triple-product identity for $\psi$); the engine's basis would need a
\verb|q_pochhammer_transform| operator (one of the four queued for
v0.10; see \cref{sec:roadmap}) to discover this rewrite autonomously
rather than through preprocessing.

\subsection{Numerical verification}\label{sec:ab166-verification}

Independently reproduced in our local environment by
\verb|analysis/29_warnaar_17_ab_166_cost_verification.py|
\cite[diary~43]{ProjectC}. Numerical agreement CONFIRMED; cost-ratio
claims \textbf{CONFIRMED-with-caveat} at $\beta = 10^{-3}$ (form D
suffers catastrophic cancellation at $\dps = 80$; see
\cref{sec:conditioning}), CONFIRMED at $\beta = 10^{-2}$ vs LHS,
PARTIAL at $\beta = 10^{-1}$ vs product spelling.

Identity values (all forms with valid numerics agree to $\ge 72$
digits at the $\beta = 0.1$ cross-check):
\begin{center}
\begin{tabular}{r l r}
\toprule
$\beta$ & $q = e^{-\beta}$ & $\psi(e^{-\beta})$ value \\
\midrule
$0.001$ & $0.99900\ldots$ & $39.638$ \\
$0.01$  & $0.99004\ldots$ & $12.549$ \\
$0.1$   & $0.90483\ldots$ & $4.013$  \\
\bottomrule
\end{tabular}
\end{center}

Wall-clock comparison (median of $5$ trials, \texttt{mpmath} $80$
$\dps$; full table in
\verb|analysis/results/29_warnaar_17_ab_166_cost_verification.csv|):

\begin{center}
\small
\begin{tabular}{r r r r r r r l}
\toprule
$\beta$ & $L$ tri.\ (D) & $\psi$ prod.\ (E) & Project~C (F) & F vs D & F vs E & claim (vs D) & verdict \\
\midrule
$0.001$ & $28.6$ ms*  & $1{,}079$ ms & $0.0333$ ms & $\mathbf{858\times^{*}}$ & $\mathbf{32{,}363\times}$ & $34\times$ & CONFIRMED ($\ge 25\times$) \\
$0.01$  & $3.85$ ms   & $10.4$ ms    & $0.0334$ ms & $\mathbf{118\times}$     & $\mathbf{317\times}$      & $15\times$ & CONFIRMED ($\ge 8\times$)  \\
$0.1$   & $0.839$ ms  & $1.89$ ms    & $0.0309$ ms & $\mathbf{26.7\times}$    & $\mathbf{60\times}$       & $9\times$  & CONFIRMED ($\ge 3\times$)  \\
\bottomrule
\end{tabular}
\end{center}

\noindent
$^{*}$Form D at $\beta = 10^{-3}$ returns $-9.76 \times 10^{96}$
instead of the true value $+39.638$---catastrophic cancellation
across $10^{177}$ intermediate magnitudes at $\dps = 80$. The
$858\times$ ratio measures the work the naive evaluator does, not the
work to actually recover $30$ digits. For $30$-digit-correct recovery
via form D, $\dps \gtrsim 220$ is needed; form F (Project~C) recovers
$30$ digits at $\dps = 80$ in $\sim 33\,\mu$s.

\textbf{Speedups.} $27$--$858\times$ vs the LHS triangular sum
(caveat at $\beta = 10^{-3}$); $60$--$32{,}363\times$ vs the product
spelling near $q \to 1^-$. The claimed direction in
\cite[diary~42]{ProjectC} ($9$--$34\times$ vs LHS; ``hundreds to
tens of thousands'' vs product) is confirmed; the measured
\emph{magnitude} of the LHS-comparison speedup exceeds the claim
because form D at $\beta = 10^{-3}$ is much slower than the
ChatGPT-Pro measurement reported (likely platform / version
difference).

\subsection{A secondary point: numerical conditioning}\label{sec:conditioning}

At $\beta = 10^{-3}$, the LHS hypergeometric summands reach
magnitudes $\sim 10^{177}$ before cancelling to the true value
$\sim 40$. Naive $30$-digit computation of the LHS is \emph{not just
slow, it is numerically ill-conditioned}---large guard precision is
required to recover the small cancelled value. The transformed form
has no such issue: every term in the dual sum is bounded by $1$.

This is an \emph{additional} practical benefit of the
transformation, on top of the wall-clock speedup: the engine is
producing a form that is better-conditioned, not merely faster.

\subsection{Discussion}

This identity is more striking than Result~1 in one respect (the
$\sim 10^{177}$ cancellation hazard) and less striking in another
(the engine's basis cannot currently discover the Jacobi
triple-product rewrite that takes $L(q)$ to the half-shifted Gaussian
---that is a manual preprocessing step). Both shortcomings are
addressed in v0.10 (see \cref{sec:roadmap}).

\section{Methodological note: the generalized obstruction principle}\label{sec:obstruction}

A central methodological pillar---closed at high confidence over
three independent instances during the run
\cite[diaries~19, 35, 37]{ProjectC}---is the following:

\begin{principle}[Generalized obstruction principle]\label{prin:obstruction}
Scalar time-only functional calculus of a Fokker--Planck-type
generator $L$ preserves $\ker L$ (the equilibrium sector).
Therefore, any Mellin-pair-style reciprocal-symmetric spectral
object (Watson's $M(\alpha)$, Berndt--Dixit's
$BD_{\mathrm{int}}(\alpha)$, etc.) cannot be matched by a functional
of $L$'s equilibrium partition function under such a calculus.
\end{principle}

The three confirming instances:
\begin{enumerate}
\item \textbf{B1}---modular S-transform on a lattice-gas $g^*$ vs
  Watson's $M(\alpha)$: closed via $20$ candidate shadow functions
  across $5$ structurally-different families all failing the strict
  $\le 1\%$ / $\ge 2$-dex bar \cite[diary~19]{ProjectC};
  cohomological irreducibility argument. The shadow $g^*$ is an
  Eichler cocycle in $H^1(\Gamma, V_k)$.
\item \textbf{B1.5}---Caputo time-fractional Fokker--Planck with
  harmonic confinement vs Berndt--Dixit
  $BD_{\mathrm{int}}(\alpha)$: closed via a one-line argument
  (Mittag-Leffler $E_\sigma(0) = 1$ implies the equilibrium density
  is fixed pointwise under the subordination, so the partition
  function is $\sigma$-independent \cite[diary~35]{ProjectC}).
\item \textbf{B1.5-ext}---Riemann--Liouville time-fractional
  Fokker--Planck with \emph{quartic} confinement vs Berndt--Dixit:
  closed via the same one-line argument with explicit handling of
  the RL initialization source term \cite[diary~37]{ProjectC}. This
  sharpens the scope: the principle requires \emph{scalar time-only}
  functional calculus. \emph{Space-dependent} variable-order
  subdiffusion (Fedotov--Falconer literature) is the identified
  route to a potential counterexample.
\end{enumerate}

\subsection{Implication for Project~C}

The obstruction principle directly constrains Project~C's search:
\emph{any operator sequence that attempts to bridge a partition-function-
of-a-generator object on one side with a Mellin-pair-reciprocal
spectral object on the other side is, generically, structurally
ineligible.} The engine's basis does not include such bridge attempts
(operators 1--3 all act \emph{within} the analytic or modular
sector, not across them); the principle gives a \emph{justification}
for the restriction, not just an empirical observation. A future
\verb|SectorClassifier| pre-pass (queued for v1.x) would make this
restriction enforceable as a hard prune step (per
\cref{sec:toolkitsplit-implication}).

\section{Connection to the toolkit-split cog-sci finding}\label{sec:toolkitsplit}

The operator basis was distilled from empirical hand-tagging of
Ramanujan's recorded work against a working vocabulary of operators.
The current state (as of \cite[diary~36]{ProjectC}; $200+$ tagged
proof units across four chapters of Berndt's commentaries and the
Quarterly Reports) supports a \emph{hierarchical-modular}
decomposition (per \cite[\S D1]{ProjectC}):

\begin{center}
\small
\begin{tabular}{l l l}
\toprule
Level & Cluster & Example chapters \\
\midrule
Top            & modular $\perp$ analytic                       & zero cross-firing in $\sim 200$ units \\
Sub (modular)  & $\{$q-shift, modular-S$\}$ for CFs             & NB V Ch.~32 \\
Sub (modular)  & $\{$modular-S$\}$ monoculture for modular eqs. & NB V Ch.~36 \\
Sub (modular)  & $\{$class-field, modular-S$\}$ for class inv.  & NB V Ch.~34 \\
Sub (analytic) & Bipartite: Mellin/Master-Theorem $\perp$ Schr\"oder--K\"onig & NB Part~1 Ch.~3--4, QRs \\
\bottomrule
\end{tabular}
\end{center}

The bipartite analytic finding \cite[diary~36;
BIG\_IDEAS\_REVIEW \S 2c]{ProjectC} is that within the analytic
toolkit, Mellin/Master-Theorem material and Schr\"oder--K\"onig
fractional iteration material fire in \emph{disjoint entry-sets}---
zero co-tagging across the $52$ proof units that constitute the
iteration and Master-Theorem subdomains of NB Part~1 Ch.~4.

\subsection{Implication for the engine}\label{sec:toolkitsplit-implication}

The basis described in \cref{sec:engine} is the \emph{engineering
image} of this empirical decomposition: the three
substantively-implemented operators correspond to the three
highest-frequency operator clusters in the tagged corpus (modular
family $\to$ operator~3; Mellin/Master-Theorem $\to$ operators~1--2;
the four stub operators are placeholders for the Schr\"oder--K\"onig
sub-cluster, Frullani, Fourier self-reciprocity, and Mellin-pair
reciprocity, which would be wired in v1.x).

The two results in \cref{sec:result1,sec:result2} both use
operator~3 (the modular family). This matches
\cite[diary~36]{ProjectC}'s finding that Lost Notebook entries in
the partial-theta / theta-function neighbourhood (Warnaar's
identity, $\psi(q)$) are exactly the targets a modular-cluster
operator should hit.

A \verb|SectorClassifier| pre-pass (queued v1.x, per
\cite[diary~38]{ProjectC}) would pin each new search state to either
the modular or analytic sector up-front and constrain the proposer
to operators from that sector only. By the toolkit-split hypothesis
(zero cross-firing observed), this is expected to give a
$\ge 3\times$ reduction in proposer-space search effort with no loss
of coverage.

\section{Discussion and future work}\label{sec:discussion}

\subsection{What was demonstrated}

This paper reports \emph{two specific Ramanujan identities where an
automated symbolic-search engine, equipped with a small operator
basis distilled from Ramanujan's empirically-observed working
clusters, finds transformed computational forms that are
$13.7\times$ to $32{,}363\times$ faster than the recorded forms at
$q$ near $1^-$}. Both transformations use only the engine's existing
\verb|modular_transform_with_correction| operator (Result~1) and
\verb|mellin_inversion| operator (Result~2); both are verified to
$\le 6 \cdot 10^{-72}$ relative mismatch at $80$-digit precision at
the cross-check cell $\beta = 0.1$.

\subsection{What this is NOT}

It is important to be explicit:
\begin{enumerate}
\item The identities are \emph{recorded by Ramanujan} (Lost Notebook
  pp.~29 and 35) and have been \emph{verified by Andrews, Berndt,
  Warnaar and others} in the standard commentaries. The engine did
  \emph{not} discover them.
\item The \emph{transformed forms} are also classical in the sense
  that the underlying modular transformations (Dedekind-$\eta$,
  $\theta_2 \leftrightarrow \theta_4$) are textbook. The engine's
  contribution is to \emph{compose them with the right
  preprocessing, verify numerically to high precision, and pass a
  cost gate that confirms the order-of-magnitude speedup is real}.
\item The cost-ratio numbers in \cref{sec:result1,sec:result2} have
  been \emph{independently reproduced} in our \texttt{mpmath}
  environment \cite[diary~43]{ProjectC} and are within a factor of
  $2\times$ of the originally claimed ChatGPT-Pro measurements
  \cite[diary~42]{ProjectC} at every $\beta$ cell.
\item Both routes use \emph{algebraic preprocessing} (parameter
  binding, Pochhammer rewrite for \cref{sec:result1}; Jacobi
  triple-product rewrite for \cref{sec:result2}) that is currently
  done by hand. \Cref{sec:ab166-verification,sec:roadmap} discuss
  the v0.10 operator additions that would close this gap.
\end{enumerate}

\subsection{Basis-completeness and the v0.10 roadmap}\label{sec:roadmap}

As of v0.9, Project~C reaches $\mathbf{8}$ of $\mathbf{12}$ seed
targets at depth $\le 4$ under a budget of $10{,}000$. The four
unreachable targets---mock-theta $f_3$, partial-theta sums, finite
Gauss-sum evaluations, and generic $q$-Pochhammer transforms---have
no operator whose \verb|apply_fn| currently returns non-\verb|None|
on them. v0.10 adds four operators:
\begin{enumerate}
\item \texttt{mock\_theta\_completion\_transform}: thin wrapper
  around the Zwegers $R_{a,b}$ erfc form
  \cite[diary~41]{ProjectC}; see also \cite{Zwegers2002}. Resolves
  a known methodological bug in our P2 prompt (which used a
  divergent integral form; the convergent Mordell form in
  \verb|analysis/17_mock_theta_residual_ebm_v1.py| is what our code
  actually uses).
\item \texttt{partial\_theta\_completion\_or\_false\_theta\_transform}:
  completes partial-theta sums $\sum q^{n(n+1)/2}$ (optionally with
  $(-1)^n$).
\item \texttt{gauss\_sum\_finite\_evaluation}: cleanest of the four;
  quadratic Gauss-sum closed form with the silent-zero guard
  preserved (per the documented hazard at $q \equiv 2 \pmod 4$ in
  \verb|rkc/registry.py|).
\item \texttt{q\_pochhammer\_transform}: most general; lifts seed
  factors to recognised infinite-product form and applies
  Dedekind-$\eta$ or Heine / $q$-binomial as appropriate. Would
  close the AB~1.6.6 preprocessing gap noted in
  \cref{sec:ab166-verification}.
\end{enumerate}

ChatGPT~Pro estimates $1$--$2$ days each
\cite[diary~42]{ProjectC}; with realistic test coverage and edge
cases, more likely $3$--$7$ days each. After v0.10: Project~C would
close the basis-completeness milestone for all $12$ seed targets and
unblock $4$ registry families currently surfacing as
\verb|CompletionNotImplemented| in \verb|rkc/completions.py|.

\subsection{The from-scratch-discovery milestone}

The natural next milestone---and one that this paper does \emph{not}
yet clear---is for the engine to produce an identity that is \emph{not
recorded anywhere in Ramanujan's notebooks or Berndt's
commentaries}, and that holds to the same strict $\le 10^{-30}$
acceptance bar across a $\ge 2$-dex sweep. To our knowledge no
automated identity-discovery engine has cleared this bar for
$q$-series identities at the Ramanujan level of structure.
Project~C's roadmap to that milestone:
\begin{enumerate}[label=(\alph*)]
\item Wire the v0.10 four operators above.
\item Add the \verb|SectorClassifier| pre-pass per
  \cref{sec:toolkitsplit-implication}.
\item Add a \verb|CompressionFilter| per the metric in
  \cref{sec:intro} of this note (and \cite[diary~39]{ProjectC}),
  with the grammar-relative caveats explicit.
\item Run a directed sweep against seeds in the Lost Notebook
  neighbourhood (Warnaar~1.7 family, mock-theta $f$-family, etc.)
  at depth $\le 6$, budget $\sim 10^5$, with the engine's existing
  cost-gate set to require ratio $> 10\times$.
\end{enumerate}

Whether (a)--(d) suffice for a from-scratch hit is the open question.

\section*{Acknowledgments}

The eight ChatGPT-Pro Erd\H{o}s-style prompt cycles
\cite[diaries~35, 37, 38, 39, 40, 41, 42, plus the v0.8a
proposer-test that preceded them]{ProjectC} produced the structural
insights and operator-completion proposals that drove the
implementation reported here; their methodology (open-ended T/F
prompts with anti-handwave guards and explicit numerical
verification bars) is itself a piece of the methodology
contribution. The standing-on-shoulders base is the immense body of
work by B.C.\ Berndt, G.E.\ Andrews, S.O.\ Warnaar, F.\ Garvan,
J.\ Frye, S.\ Zwegers, D.\ Hickerson, B.C.\ Berndt and A.\ Dixit,
and the OEIS, without which none of the source identities,
transformations, or verification points would exist.
\cite[diary~41]{ProjectC}'s detection of an error in our own P2
prompt (a divergent integral form picked up from training-memory
rather than copied from our implementation) is a worked example of
the ``paste actual implementation, not memory-paraphrase''
discipline that the ChatGPT-Pro methodology imposes.

\appendix

\section{The operator-3 apply-fn (excerpted from \texttt{rkc/search/basis.py})}\label{app:apply-fn}

The pattern matchers for the two cases used in
\cref{sec:result1,sec:result2} are reproduced here in skeletal form
(full implementation in \verb|rkc/search/basis.py| at function
\verb|_modular_transform_with_correction_apply|):

\begin{verbatim}
def _modular_transform_with_correction_apply(expr, n):
    # CASE A: log Bose-Einstein  log(1 - exp(-beta n))
    #   -> log(1 - exp(-4 pi^2 n / beta))
    #      + correction (-pi^2/(6 beta) - (1/2) log(beta/(2 pi)) + beta/24)
    result = _log_bose_einstein_apply(expr, n)
    if result is not None:
        return result
    # CASE B: Lambert / E_2 quasi-modular  n q^n / (1 - q^n)
    #   -> -(4 pi^2 / beta^2) n / (exp(4 pi^2 n / beta) - 1)
    #      + correction pi^2/(6 beta^2) + 1/24 - 1/(2 beta)
    result = _lambert_e2_apply(expr, n)
    if result is not None:
        return result
    # CASE C: Stieltjes occupation-number (v0.8c)  1/(exp(beta n) - 1)
    result = _v08c_try_occupation_number_pattern(expr, n)
    if result is not None:
        return result
    return None
\end{verbatim}

Each branch returns either \verb|None| (declines) or a $3$-tuple
\verb|(new_summand, additive_correction, summand_range)|. The engine
populates the \verb|OperatorApplication| from the tuple; the
verifier picks up the correction and folds it into the before-side
residual before comparison.

% Include all entries from refs.bib in the printed bibliography,
% so that the standard Ramanujan/Berndt/Andrews-Berndt literature
% inherited from the source markdown's References block prints even
% if not directly \cite'd in the body.
\nocite{AB2009,AB2012,AB2013,BerndtNB1,BerndtNB5,BerndtDixit2009,DixitBerndt1995,BerndtBhargava2022,BerndtYee2002,Hickerson1988,Ono2015}

\bibliographystyle{amsalpha}
\bibliography{refs}

\end{document}
